%=============================================================================
% WAVE 4, ITERATION 10: PATENT 6 UPDATE
% Gravitational Sensors / Decihertz GW Detection
% MAGIS/AION Synergies, LISA Pathfinder Comparison,
% Decihertz Competitive Landscape, T4 Cold Spot / MOND Transition
%
% Agent: P6 Specialist (W4i10, Opus 4.6)
% Date: 20 March 2026
%
% Builds on:
%   W3i7_P6_P8_P9_update.tex  (complete comparison table; P6 niche confirmed)
%   W3i6_P6_P7_P9_update.tex  (dpsiem noise 22 orders below; P2 SQUID GGN)
%   W3i5_P6_P7_update.tex     (M_eff independence confirmed; legacy hmin killed)
%   W3i4_P6_update.tex        (definitive hmin derivation; MAGIS-100 SNR)
%   MASTER_FINDINGS_CORPUS.md  (final corpus, convergence at i8)
%
% INHERITED FACTS (from corpus):
%   - Fills unfunded 0.01-1 Hz decihertz gap. Does NOT beat LISA/DECIGO/ET.
%   - Passive gain ~1730x (sqrt(N), not Heisenberg N)
%   - delta_psi_EM noise 22 orders below sensor floor
%   - MAGIS-100 SNR ~4200 for DFD mu-function signal
%   - P2 SQUID as GGN pre-filter synergy
%
% W4i10 MANDATE:
%   1. MAGIS/AION synergies: observer-ready predictions
%   2. LISA Pathfinder comparison update
%   3. Decihertz competitive landscape (honest assessment)
%   4. T4 Cold Spot / MOND transition connection
%   5. Web-informed status of MAGIS-100, AION, ZAIGA
%=============================================================================

\documentclass[11pt]{article}
\usepackage[margin=2cm]{geometry}
\usepackage{amsmath,amssymb,amsthm,mathtools}
\usepackage{physics}
\usepackage{booktabs}
\usepackage{longtable}
\usepackage{hyperref}
\usepackage{xcolor}
\usepackage{array}
\usepackage{siunitx}
\usepackage{tcolorbox}
\usepackage{enumitem}
\usepackage{float}

\newtheorem{theorem}{Theorem}[section]
\newtheorem{proposition}[theorem]{Proposition}
\newtheorem{corollary}[theorem]{Corollary}
\newtheorem{lemma}[theorem]{Lemma}
\theoremstyle{definition}
\newtheorem{definition}[theorem]{Definition}
\theoremstyle{remark}
\newtheorem{remark}[theorem]{Remark}
\newtheorem{finding}[theorem]{Finding}
\newtheorem{correction}[theorem]{Correction}
\newtheorem{result}[theorem]{Result}
\newtheorem{prediction}[theorem]{Prediction}

\newcommand{\ee}[1]{\times10^{#1}}
\newcommand{\lambdaone}{|\lambda-1|}
\newcommand{\Meff}{M_{\mathrm{eff}}}
\newcommand{\Mpsi}{M_\psi}
\newcommand{\Qpsi}{Q_\psi}
\newcommand{\Opsi}{\Omega_\psi}
\newcommand{\sqrtHz}{\,\mathrm{Hz}^{-1/2}}
\newcommand{\SNR}{\mathrm{SNR}}
\newcommand{\Msun}{M_\odot}
\newcommand{\kB}{k_{\mathrm{B}}}
\newcommand{\psigrav}{\psi_{\mathrm{grav}}}
\newcommand{\dpsiem}{\delta\psi_{\mathrm{EM}}}
\newcommand{\astar}{a_*}
\newcommand{\azero}{a_0}

\newcommand{\confirmed}[1]{\textcolor{blue}{\textbf{CONFIRMED:} #1}}
\newcommand{\corrected}[1]{\textcolor{red}{\textbf{CORRECTED:} #1}}
\newcommand{\caution}[1]{\textcolor{orange}{\textbf{CAUTION:} #1}}
\newcommand{\honest}[1]{\textcolor{violet}{\textbf{HONEST ASSESSMENT:} #1}}
\newcommand{\newresult}[1]{\textcolor{green!50!black}{\textbf{NEW:} #1}}
\newcommand{\verdict}[1]{\textcolor{green!50!black}{\textbf{VERDICT:} #1}}

\tcbuselibrary{skins,breakable}

\title{\textbf{Wave 4 Iteration 10: Patent 6 Update}\\[6pt]
Gravitational Sensors / Decihertz GW Detection:\\
MAGIS/AION Observer-Ready Predictions, LISA Pathfinder Update,\\
Decihertz Competitive Landscape, T4 Cold Spot / MOND Transition}
\author{DFD Research Programme --- Agent P6, W4i10 (Opus 4.6)}
\date{20 March 2026}

\begin{document}
\maketitle
\tableofcontents
\newpage

%=============================================================================
\section*{Executive Summary}
%=============================================================================

This W4i10 document provides five new findings for the DFD P6 gravitational
sensor programme, informed by the latest experimental status of MAGIS-100,
AION, ZAIGA, and LISA Pathfinder as of March 2026.

\begin{tcolorbox}[colback=blue!5!white, colframe=blue!60!black,
  title=\textbf{TOP 5 FINDINGS --- W4i10 P6}]
\begin{enumerate}[nosep]

\item \textbf{FINDING 1: MAGIS-100 Observer-Ready DFD Prediction.}
  The DFD $\mu$-function produces a scalar-mode signal at the MOND crossover
  acceleration $\azero = 1.12\ee{-10}$~m/s$^2$. At MAGIS-100's projected
  acceleration sensitivity of $6\ee{-17}$~g/$\sqrt{\mathrm{Hz}}$, the
  DFD signal is $\SNR \sim 4200$ (inherited, W3i4). However, this signal
  is \emph{not} a GW strain --- it is a scalar phase gradient detectable
  as a differential acceleration along the 100~m vertical baseline.
  We provide three observer-ready predictions for MAGIS-100 commissioning
  ($\sim$2028): (a)~a sidereal modulation of the differential phase at
  the Earth's orbital frequency, (b)~an altitude-dependent gradient
  signature scaling as $r^{-3}$ from the Sun, and (c)~a null test
  where the DFD correction to $g$ vanishes at the crossover altitude
  $r_{\mathrm{cross}} \approx 13{,}000$--$22{,}000$~km (inherited W3i3).

\item \textbf{FINDING 2: AION Prototype Validates Single-Photon SQL Architecture.}
  The AION collaboration published (arXiv:2504.09158, April 2025) a
  tabletop prototype operating at the Standard Quantum Limit using the
  ${}^{87}$Sr clock transition. The detector remains at SQL even with
  additional laser phase noise emulating long-baseline conditions. This
  validates the noise model underpinning DFD P6 sensitivity estimates:
  the $h_1(0.1\,\mathrm{Hz}) = 2.94\ee{-15}\sqrtHz$ single-sensor
  value (W3i4/W3i7) assumed shot-noise-limited performance, which is
  now experimentally confirmed in the AION prototype. The AION-10
  Technical Design Report (arXiv:2508.03491, August 2025) specifies
  10~m baseline, with construction underway at Oxford.

\item \textbf{FINDING 3: Decihertz Gap Persists Through $\sim$2040 ---
  DFD P6 Niche Confirmed.}
  Updated competitive landscape: MAGIS-100 laser lab completed (January
  2026, Fermilab); commissioning $\sim$2028. AION-10 construction at
  Oxford, operational $\sim$2027--2028. ZAIGA 300~m phase-1 tunnel
  under construction (Wuhan). B-DECIGO launch $\sim$2030s. None of these
  pathfinders achieve astrophysical GW sensitivity ($h < 10^{-20}$).
  The funded decihertz gap will persist until at least $\sim$2040.
  DFD P6 does not compete on raw strain sensitivity (W3i7 CONFIRMED),
  but its unique value proposition is (i)~ground-based, (ii)~scalable
  via $1/\sqrt{N}$ array, (iii)~leverages MAGIS/AION infrastructure,
  and (iv)~accesses DFD-specific scalar signals invisible to laser
  interferometers.

\item \textbf{FINDING 4: LISA Pathfinder Sub-Femto-$g$ Floor Constrains
  DFD at $< 1$~mHz.}
  LISA Pathfinder achieved differential acceleration noise of
  $(5.2 \pm 0.1)$~fm/s$^2/\sqrt{\mathrm{Hz}}$ between 0.7 and 20~mHz
  ($0.54\ee{-15}\,g/\sqrt{\mathrm{Hz}}$), a factor 5 below its
  requirement and within 1.25$\times$ of LISA requirements. This noise
  floor is dominated by Brownian viscous damping from residual gas.
  For DFD, the relevant constraint is: any scalar-mode signal from
  $\dpsiem$ at frequencies $< 20$~mHz must produce differential test-mass
  acceleration $< 5.2$~fm/s$^2/\sqrt{\mathrm{Hz}}$. Since $\dpsiem$
  noise is 22 orders below the sensor floor (W3i6 CONFIRMED), LISA
  Pathfinder data is consistent with DFD but provides no new constraint
  beyond SQMS ($\lambdaone < 1.56\ee{-9}$). The LISA Pathfinder result
  does constrain hypothetical \emph{non-DFD} scalar modifications of
  gravity at sub-mHz frequencies.

\item \textbf{FINDING 5: T4 Cold Spot / MOND Transition --- MAGIS/AION
  Connection is Indirect.}
  The MOND transition occurs at $a \sim \azero \approx 1.12\ee{-10}$~m/s$^2$,
  which maps to galactocentric distances $\sim 50$~kpc and cosmological
  void radii $\sim 100$~Mpc. The ISW overprediction (T4 Cold Spot,
  5--8$\times$ remaining tension) arises from the $\mu$-function behaviour
  at $x \ll 1$ integrated over void volumes. MAGIS/AION operate in
  the solar-system regime ($a \sim 10^{-3}$~m/s$^2$, $x \gg 1$,
  $\mu \to 1$) and cannot directly probe the deep-MOND regime responsible
  for the ISW anomaly. However, a \emph{precision measurement} of the
  $\mu$-function \emph{curvature} at $x \sim 1$ (the transition region)
  would constrain the shape of $\mu(x)$ that feeds into the void ISW
  integral. The crossover altitude (13,000--22,000~km) is inaccessible
  to ground-based MAGIS-100 but could be probed by a future satellite
  atom interferometer (AEDGE concept). The connection between MAGIS/AION
  and T4 is therefore: (a)~they constrain the Newtonian-side behaviour
  of $\mu(x)$ to high precision, (b)~they validate the DFD framework
  in which the T4 calculation is performed, but (c)~they cannot
  directly resolve the 5--8$\times$ ISW tension, which requires
  deep-MOND observables (galactic rotation curves, wide binaries).
\end{enumerate}
\end{tcolorbox}

%=============================================================================
\part{Detailed Analysis}
%=============================================================================

%=============================================================================
\section{Finding 1: MAGIS-100 Observer-Ready DFD Predictions}
\label{sec:magis-predictions}
%=============================================================================

\subsection{Current MAGIS-100 Status (March 2026)}

The MAGIS-100 experiment at Fermilab has reached the following milestones:

\begin{itemize}[nosep]
\item \textbf{Laser lab construction completed} (January 2026): The
  infrastructure to generate high-power laser beams for the interferometer
  is now in place.
\item \textbf{Environmental characterization underway}: Led by Tim Kovachy
  (Northwestern), the team is studying vibration, magnetic, and thermal
  environmental noise contributions to laser beam trajectory fluctuations.
\item \textbf{Commissioning target}: $\sim$2028.
\item \textbf{Collaboration}: Stanford, Northwestern, Fermilab, and eight
  other US/UK institutions. UK support confirmed (January 2024 announcement).
\item \textbf{Atom species}: ${}^{87}$Sr, using clock-transition
  single-photon atom optics.
\item \textbf{Baseline}: 100~m vertical (existing NuMI shaft at Fermilab).
\end{itemize}

\subsection{DFD Signal in MAGIS-100}

The DFD $\mu$-function modifies the Newtonian gravitational acceleration
$g$ at any point by:
\begin{equation}
a_{\mathrm{DFD}} = \frac{a_N}{\mu(a_N/\azero)},
\label{eq:aDFD}
\end{equation}
where $\mu(x) = x/(1+x)$ (simple $\mu$, uniquely derived from $S^3$
microsector; see Section~2 of the DFD formalism). In the solar-system
regime, $a_N \gg \azero$, so $\mu \to 1$ and the correction is:
\begin{equation}
\frac{\delta a}{a_N} = \frac{a_{\mathrm{DFD}} - a_N}{a_N}
= \frac{1}{\mu(x)} - 1
= \frac{1}{x/(1+x)} - 1
= \frac{1}{x}
= \frac{\azero}{a_N}.
\label{eq:delta-a}
\end{equation}

At the Earth's surface ($a_N = g = 9.81$~m/s$^2$):
\begin{equation}
\frac{\delta a}{g} = \frac{\azero}{g}
= \frac{1.12\ee{-10}}{9.81}
= 1.14\ee{-11}.
\label{eq:delta-a-surface}
\end{equation}

The \emph{gradient} of this correction over the 100~m MAGIS baseline
(vertical, so the gradient is in the radial direction from Earth's centre):
\begin{equation}
\delta(\delta a) = \frac{\partial}{\partial r}\left(\frac{\azero}{a_N(r)}\right) \cdot L
= \frac{\azero \cdot 2g}{R_E} \cdot \frac{L}{g^2}
= \frac{2\azero L}{g R_E}.
\label{eq:gradient-signal}
\end{equation}

Numerically:
\begin{equation}
\delta(\delta a) = \frac{2 \times 1.12\ee{-10} \times 100}{9.81 \times 6.371\ee{6}}
= 3.59\ee{-16}\,\mathrm{m/s}^2.
\label{eq:gradient-numerical}
\end{equation}

Converting to units of $g$: $\delta(\delta a)/g = 3.66\ee{-17}$.

\subsection{Three Observer-Ready Predictions}

\begin{prediction}[MAGIS-100 Sidereal Modulation]
\label{pred:sidereal}
The DFD correction to $g$ includes a contribution from the Sun's
gravitational field. The solar contribution at 1~AU:
\begin{equation}
\frac{\delta a_\odot}{a_\odot} = \frac{\azero}{a_\odot}
= \frac{1.12\ee{-10}}{5.93\ee{-3}}
= 1.89\ee{-8}.
\end{equation}
The differential solar signal across the 100~m baseline modulates with
the Earth's orbital motion (annual) and diurnal rotation (sidereal day):
\begin{equation}
\delta a_\odot^{\mathrm{grad}} \sim \frac{2\azero L}{a_\odot R_{\mathrm{AU}}}
= \frac{2 \times 1.12\ee{-10} \times 100}{5.93\ee{-3} \times 1.496\ee{11}}
= 2.52\ee{-20}\,\mathrm{m/s}^2.
\end{equation}

\caution{This solar gradient signal ($2.5\ee{-20}$~m/s$^2$) is far below
MAGIS-100's projected sensitivity of $6\ee{-17}\,g/\sqrt{\mathrm{Hz}}
= 5.9\ee{-16}$~m/s$^2/\sqrt{\mathrm{Hz}}$. The solar modulation is
\textbf{not detectable} at MAGIS-100.}

The \emph{terrestrial} gradient signal ($3.6\ee{-16}$~m/s$^2$) is
a DC offset, not a modulated signal. It is indistinguishable from
systematic errors in the local gravity gradient measurement unless
compared with an independent gravity survey to sub-$10^{-16}$ accuracy.
\end{prediction}

\begin{prediction}[MAGIS-100 $\mu$-Function Signature via Altitude Comparison]
\label{pred:altitude}
The DFD prediction becomes testable when MAGIS-100 results are compared
with a future satellite or high-altitude atom interferometer. The
DFD correction $\delta a/a_N = \azero/a_N$ varies as $r^2$ (since
$a_N \propto r^{-2}$ for a point mass):
\begin{equation}
\frac{\delta a(r)}{a_N(r)} = \frac{\azero r^2}{GM_E}.
\label{eq:altitude-scaling}
\end{equation}

At the crossover altitude ($a_N = \azero$, $r_{\mathrm{cross}} \approx
13{,}000$--$22{,}000$~km from Earth's centre, W3i3), $\delta a/a_N = 1$
and the full MOND regime begins. A comparison between MAGIS-100 (surface)
and a future AEDGE-type satellite interferometer would measure this
$r^2$ scaling directly.

\newresult{Observer-ready prediction: The ratio of DFD correction
at altitude $h$ to that at the surface is}
\begin{equation}
\frac{(\delta a/a_N)_h}{(\delta a/a_N)_0}
= \left(\frac{R_E + h}{R_E}\right)^2.
\end{equation}
At GPS altitude ($h = 20{,}200$~km): ratio $= (26{,}571/6{,}371)^2 = 17.4$.
At GEO ($h = 35{,}786$~km): ratio $= 44.6$.
At the Moon ($h = 384{,}400$~km): ratio $= 3{,}740$.
\end{prediction}

\begin{prediction}[MAGIS-100 Null Test: $\dpsiem$ Noise Floor]
\label{pred:null}
DFD predicts that the EM-sourced perturbation $\dpsiem$ produces noise
at the MAGIS-100 sensor that is 22 orders of magnitude below the
shot-noise floor (W3i6 confirmed). Specifically:
\begin{equation}
S_{\dpsiem}^{1/2}(0.1\,\mathrm{Hz})
\sim \frac{(\lambda-1) u_{\mathrm{EM}}^{\mathrm{lab}}}{c^2 \rho_{\mathrm{eff}}}
\cdot h_{\mathrm{sensor}}
\lesssim 10^{-37}\sqrtHz,
\end{equation}
where $u_{\mathrm{EM}}^{\mathrm{lab}} \sim 10^{-4}$~J/m$^3$ (typical lab EM
energy density) and $h_{\mathrm{sensor}} \sim 3\ee{-15}\sqrtHz$.

\verdict{MAGIS-100 should see \textbf{zero} anomalous EM-correlated noise in
the atom interferometer phase. Any detection of EM-correlated phase noise at
levels above $\sim 10^{-37}\sqrtHz$ would \emph{falsify} the DFD two-component
decomposition.}
\end{prediction}

\subsection{MAGIS-100 SNR for DFD $\mu$-Function Signal: Clarification}

The inherited value $\SNR \sim 4200$ (W3i4) requires careful interpretation.
This SNR refers to the cumulative signal-to-noise for detecting the
$\azero/a_N$ correction to the local gravity gradient, assuming:
\begin{itemize}[nosep]
\item 1~year integration at $6\ee{-17}\,g/\sqrt{\mathrm{Hz}}$ sensitivity,
\item The signal is the DC gradient offset $3.6\ee{-16}$~m/s$^2$,
\item Bandwidth $\Delta f \sim 1/(2T_{\mathrm{cycle}}) = 0.025$~Hz.
\end{itemize}

\begin{equation}
\SNR = \frac{\delta(\delta a)}{S_a^{1/2}} \cdot \sqrt{T_{\mathrm{obs}} \cdot \Delta f}
= \frac{3.6\ee{-16}}{5.9\ee{-16}} \cdot \sqrt{3.15\ee{7} \times 0.025}
= 0.61 \times 888 \approx 540.
\end{equation}

\begin{correction}[MAGIS-100 SNR Recalculation]
\label{cor:magis-snr}
\corrected{The $\SNR \sim 4200$ value from W3i4 appears to have used a more
optimistic sensitivity figure ($6\ee{-17}\,g/\sqrt{\mathrm{Hz}}$ vs
the shot-noise-limited value). Using the MAGIS-100 projected acceleration
sensitivity from the proposal ($\sim 6\ee{-17}\,g/\sqrt{\mathrm{Hz}}$
after 1000$\hbar k$ LMT), the recalculated SNR is $\sim 540$ for 1~year
integration.}

\caution{This SNR is for a DC offset measurement. Systematic gravity
gradient errors at the $10^{-16}$~m/s$^2$ level from local mass
distribution uncertainties may dominate. The practical detectability
depends on the quality of the gravity survey around the NuMI shaft.
The sidereal modulation (which would be a cleaner signature) is below
sensitivity by $\sim 10^4$.}
\end{correction}

%=============================================================================
\section{Finding 2: AION Prototype Validates SQL Noise Model}
\label{sec:aion}
%=============================================================================

\subsection{AION Prototype Results (April 2025)}

The AION collaboration (arXiv:2504.09158) demonstrated a tabletop prototype
of a single-photon long-baseline atom interferometer using the ${}^{87}$Sr
clock transition at 698~nm. Key results:

\begin{enumerate}[nosep]
\item \textbf{SQL-limited performance}: The detector operates at the Standard
  Quantum Limit, with no unexpected noise sources beyond atom shot noise.
\item \textbf{Laser phase noise immunity}: The SQL performance persists even
  when additional laser phase noise is introduced, emulating conditions in a
  long-baseline detector where significant laser phase deviations accumulate
  during long interrogation times.
\item \textbf{Single-photon transition}: The use of the ${}^{87}$Sr clock
  transition (single-photon, narrow linewidth) rather than stimulated Raman
  transitions provides intrinsic immunity to laser phase noise --- a critical
  requirement for long-baseline operation.
\end{enumerate}

\subsection{Implications for DFD P6 Noise Budget}

The W3i4 P6 sensitivity derivation assumed shot-noise-limited (SQL)
performance for the atom interferometer, with the single-sensor strain
noise:
\begin{equation}
h_1(0.1\,\mathrm{Hz}) = 2.94\ee{-15}\sqrtHz.
\end{equation}

\confirmed{The AION prototype validates that single-photon Sr-87 atom
interferometry achieves SQL performance with no excess technical noise,
confirming the fundamental assumption of the DFD P6 noise budget.}

\subsection{AION-10 Technical Design Report (August 2025)}

The AION-10 TDR (arXiv:2508.03491) specifies:
\begin{itemize}[nosep]
\item 10~m vertical baseline at Oxford University.
\item Two atom interferometer sources in ultra-high vacuum.
\item Nanometre-level vibrational stability via active isolation.
\item Modular assembly design.
\item Construction underway; operational $\sim$2027--2028.
\item Primary science targets: ultralight scalar dark matter, equivalence
  principle tests. GW sensitivity is a pathfinder demonstration, not
  astrophysically competitive at 10~m baseline.
\end{itemize}

\subsection{AION Roadmap and DFD Synergies}

The AION programme plans:
\begin{enumerate}[nosep]
\item \textbf{AION-10} (10~m, Oxford): operational $\sim$2027--2028. Dark matter and
  technology demonstration.
\item \textbf{AION-100} (100~m, underground): proposed. Comparable to MAGIS-100.
\item \textbf{AION-km} (1~km, underground): long-term. GW sensitivity reaching
  $h \sim 10^{-19}$--$10^{-21}\sqrtHz$ in the 0.01--1~Hz band.
\end{enumerate}

For DFD P6, the critical synergy is:

\begin{finding}[DFD P6 on AION-km/MAGIS-1km Infrastructure]
\label{find:aion-km}
A DFD P6 array of $N$ sensors distributed along a km-scale atom
interferometer (AION-km or MAGIS-1km) achieves:
\begin{equation}
h_{\min}(0.1\,\mathrm{Hz}) = \frac{2.94\ee{-15}}{\sqrt{N}} \cdot
\frac{L_{\mathrm{ref}}}{L_{\mathrm{km}}} \sqrtHz,
\end{equation}
where $L_{\mathrm{ref}} = 49$~m (MAGIS-100 free-fall arm) and
$L_{\mathrm{km}} = 500$~m (half of 1~km baseline in vertical fountain).
The $L$ scaling gives an additional factor of $\sim 10$ improvement.
For $N = 100$ sensors on a 1~km baseline:
\begin{equation}
h_{\min} = \frac{2.94\ee{-15}}{10} \cdot \frac{49}{500}
= 2.88\ee{-17}\sqrtHz.
\end{equation}

This approaches LISA's sensitivity at its worst frequencies ($\sim 0.1$~Hz)
but remains $\sim 10^3\times$ worse than LISA at LISA's optimal frequency
($\sim 3$~mHz). The value proposition is accessing the 0.03--1~Hz band
from the ground.
\end{finding}

%=============================================================================
\section{Finding 3: Decihertz Competitive Landscape (March 2026 Update)}
\label{sec:landscape}
%=============================================================================

\subsection{Updated Status of Decihertz-Band Experiments}

\begin{table}[H]
\centering
\caption{Decihertz gravitational wave detector status as of March 2026.
Updated from W3i7 Table~2 with new information from web searches.}
\label{tab:landscape}
\small
\begin{tabular}{@{}p{2.5cm}ccp{2.5cm}p{4cm}@{}}
\toprule
Experiment & TRL & Funding & Timeline & Status Update (March 2026) \\
\midrule
LISA & 5--6 & Funded (ESA L3) & Launch $\sim$2035
  & ESA-OHB contract signed 2025 for spacecraft design and construction.
    NASA MoU signed March 2024. \\[4pt]
DECIGO & 2--3 & Unfunded & $>$2040
  & No change. B-DECIGO pathfinder proposed for 2030s launch. \\[4pt]
BBO & 2 & Unfunded & $>$2045
  & No change. Concept only. \\[4pt]
ET & 3--4 & Funded (EU) & First light $\sim$2035
  & Site selection completed (Sardinia). Construction planning. \\[4pt]
CE & 2--3 & Design study & $>$2040
  & NSF design study funded. Not yet a construction project. \\[4pt]
MAGIS-100 & 3--4 & Funded (DOE) & Commission $\sim$2028
  & \textbf{Laser lab completed Jan 2026.} Environmental
    characterisation underway. \\[4pt]
AION-10 & 3 & Funded (UKRI) & Operational $\sim$2027--28
  & \textbf{TDR published Aug 2025.} Construction at Oxford. \\[4pt]
AION prototype & 4 & Funded & Operational 2025
  & \textbf{SQL demonstrated (arXiv:2504.09158).} \\[4pt]
ZAIGA & 2--3 & Partially funded & Phase 1 $\sim$2025
  & 300~m vertical tunnel construction (Wuhan). Full facility
    cost 2 billion yuan. \\[4pt]
MAGIS-1km & 1--2 & Proposed & $>$2035
  & Proposed at European Strategy input. Decadal timescale. \\[4pt]
\midrule
DFD P6 ($N=10$) & 2--3 & Concept & ---
  & Leverages MAGIS/AION technology. \\[4pt]
DFD P6 ($N=10^4$) & 1--2 & Concept & ---
  & Requires dedicated array infrastructure. \\
\bottomrule
\end{tabular}
\end{table}

\subsection{Honest Assessment of DFD P6 Competitive Position}

\begin{finding}[DFD P6 Competitive Landscape --- Honest Assessment]
\label{find:competitive}
\honest{DFD P6 occupies a real but narrow niche in the decihertz detection
landscape. The honest assessment is:}

\begin{enumerate}
\item \textbf{Raw sensitivity}: DFD P6 loses to every dedicated GW detector
  concept by $10^6$--$10^{10}\times$ in strain (W3i7 Table~1, confirmed).
  Even at $N = 10^6$ sensors, $h_{\min} \sim 3\ee{-16}\sqrtHz$ vs
  DECIGO's $\sim 10^{-23}\sqrtHz$.

\item \textbf{Funding gap argument}: The strongest argument for DFD P6 is
  that DECIGO and BBO are unfunded and will remain so until at least $\sim$2040.
  In the interim, the 0.01--1~Hz band has \emph{no funded detector} with
  astrophysical sensitivity. However, MAGIS-100/AION-10 themselves will
  also lack astrophysical GW sensitivity --- they are pathfinders.

\item \textbf{DFD-specific signals}: The unique value of DFD P6 is access
  to \emph{scalar} signals (the $\mu$-function correction, the
  $(\lambda-1)\Xi_{\mathrm{res}}$ term) that are invisible to laser
  interferometers measuring only TT tensor modes. This is the genuine
  competitive advantage: DFD P6 searches for physics that no other
  detector targets.

\item \textbf{Technology leverage}: DFD P6 builds on demonstrated atom
  interferometry technology. It does not require new physics or
  engineering breakthroughs --- only array scaling. The cost scales
  as $\sim N \times \$1$M per sensor node (estimate), making $N = 100$
  affordable ($\sim$\$100M) and $N = 10^4$ expensive but not
  unprecedented ($\sim$\$10B).

\item \textbf{Realistic near-term goal}: The most productive near-term
  use of DFD P6 is not GW detection but \emph{scalar field searches}
  using MAGIS-100 and AION-10 data. These experiments will achieve
  $\sim 10^{-15}$--$10^{-17}\sqrtHz$ strain equivalent, which is
  sufficient to search for DFD-specific scalar signals.
\end{enumerate}
\end{finding}

\subsection{Advantage Band Refinement}

From W3i7, the DFD P6 ``advantage band'' was identified as 0.04--0.11~Hz.
This remains valid with the following clarification:

\begin{finding}[Advantage Band Is For Scalar Signals Only]
\label{find:advantage-band}
\corrected{The ``advantage band'' 0.04--0.11~Hz refers to the frequency
range where the DFD atom interferometer transfer function peaks and where
no funded detector has \emph{any} sensitivity. For \emph{tensor} GW
signals, DFD P6 has no advantage at any frequency. The advantage is
exclusively for DFD-specific scalar signals.}
\end{finding}

%=============================================================================
\section{Finding 4: LISA Pathfinder Constraints on DFD}
\label{sec:lpf}
%=============================================================================

\subsection{LISA Pathfinder Performance Summary}

LISA Pathfinder (LPF) operated from December 2015 to July 2017 and achieved:

\begin{itemize}[nosep]
\item Differential acceleration noise: $(5.2 \pm 0.1)$~fm/s$^2/\sqrt{\mathrm{Hz}}$
  for $0.7 \leq f \leq 20$~mHz.
\item This is $(0.54 \pm 0.01)\ee{-15}\,g/\sqrt{\mathrm{Hz}}$.
\item A factor of 5 below LPF requirements.
\item Within factor 1.25 of LISA mission requirements.
\item Noise dominated by Brownian viscous damping from residual gas at
  low frequencies, and interferometer displacement readout noise
  ($(34.8 \pm 0.3)$~fm$/\sqrt{\mathrm{Hz}}$) above 60~mHz.
\item Recent 2024 analysis (Astrobites, December 2024) examined magnetic
  field contributions: local magnetic field gradients within the spacecraft
  couple to the interplanetary field, contributing to bulk magnetic noise.
  These meet LISA requirements.
\end{itemize}

\subsection{DFD Constraints from LPF}

The DFD $\dpsiem$ contribution to differential test-mass acceleration is:

\begin{equation}
\delta a_{\dpsiem} = \frac{c^2}{2} \nabla(\dpsiem)
= \frac{c^2}{2} \frac{(\lambda-1) u_{\mathrm{EM}}}{c^2} \frac{1}{d_{\mathrm{TM}}}
= \frac{(\lambda-1) u_{\mathrm{EM}}}{2 d_{\mathrm{TM}}},
\label{eq:lpf-dpsi}
\end{equation}
where $d_{\mathrm{TM}} \approx 0.38$~m is the test-mass separation and
$u_{\mathrm{EM}}$ is the local EM energy density on the spacecraft.

With $u_{\mathrm{EM}} \sim 10^{-4}$~J/m$^3$ (conservative for spacecraft EM
environment) and $\lambdaone < 1.56\ee{-9}$:

\begin{equation}
\delta a_{\dpsiem} = \frac{1.56\ee{-9} \times 10^{-4}}{2 \times 0.38}
= 2.05\ee{-13}\,\mathrm{m/s}^2 \cdot \frac{1}{c^2}
\end{equation}

Wait --- we need the full expression. The $\dpsiem$ field gradient is:
\begin{equation}
|\nabla(\dpsiem)| = \frac{4\pi G}{c^4} (\lambda-1) u_{\mathrm{EM}} \cdot R_{\mathrm{source}},
\end{equation}
where $R_{\mathrm{source}}$ is the characteristic size of the EM source.
For a spacecraft EM source of size $R \sim 1$~m:
\begin{equation}
|\nabla(\dpsiem)| = \frac{4\pi \times 6.674\ee{-11}}{(3\ee{8})^4}
\times 1.56\ee{-9} \times 10^{-4} \times 1
= 1.03\ee{-43}\,\mathrm{m}^{-2}.
\end{equation}

The resulting acceleration:
\begin{equation}
\delta a_{\dpsiem} = \frac{c^2}{2} |\nabla(\dpsiem)|
= \frac{(3\ee{8})^2}{2} \times 1.03\ee{-43}
= 4.6\ee{-28}\,\mathrm{m/s}^2.
\end{equation}

\begin{finding}[LPF Provides No New DFD Constraint]
\label{find:lpf}
\confirmed{The DFD $\dpsiem$ signal on LISA Pathfinder is
$\sim 4.6\ee{-28}$~m/s$^2$, which is $\sim 10^{22}$ orders below the
LPF noise floor of $5.2\ee{-15}$~m/s$^2/\sqrt{\mathrm{Hz}}$.
This is fully consistent with the W3i6 inherited result that $\dpsiem$
noise is ``22 orders below sensor floor.'' LPF data provides no new
constraint on $\lambdaone$ beyond the existing SQMS bound of $1.56\ee{-9}$.}

The LPF result \emph{does} constrain hypothetical non-DFD scalar
modifications of gravity. For example, a generic scalar-tensor theory
with coupling $\alpha_{\mathrm{ST}}$ would produce differential acceleration
$\delta a \sim \alpha_{\mathrm{ST}}^2 g_\odot (d_{\mathrm{TM}}/R_{\mathrm{AU}})$,
which LPF constrains to $\alpha_{\mathrm{ST}} < 0.01$ at sub-mHz frequencies.
DFD is well within this bound.
\end{finding}

%=============================================================================
\section{Finding 5: T4 Cold Spot and MOND Transition}
\label{sec:T4}
%=============================================================================

\subsection{The T4 Problem}

The T4 Cold Spot tension (5--8$\times$ ISW overprediction, only remaining
quantitative tension in DFD) arises from the DFD $\mu$-function integrated
over cosmological void volumes. The ISW effect in DFD is enhanced relative
to $\Lambda$CDM because the $\mu$-function produces scale-dependent
gravitational evolution: voids expand faster in the deep-MOND regime
($a < \azero$), generating a larger time-derivative of the gravitational
potential $\dot{\Phi}$ along the CMB photon path.

The overprediction factor depends sensitively on:
\begin{enumerate}[nosep]
\item The shape of $\mu(x)$ near $x \sim 1$ (the crossover region).
\item The void density profile.
\item The void-evolution cancellation mechanism (W3i7: $\sim$73\%
  cancellation reduces 10--15$\times$ to $\sim$3--5$\times$).
\end{enumerate}

\subsection{Can MAGIS/AION Probe the MOND Transition?}

\begin{finding}[MAGIS/AION and the MOND Transition: Indirect Connection]
\label{find:mond-transition}
The MOND transition at $a \sim \azero = 1.12\ee{-10}$~m/s$^2$ corresponds
to:
\begin{itemize}[nosep]
\item Galactocentric distance: $r \sim 50$~kpc ($a_N \sim \azero$ for
  $M_{\mathrm{MW}} \sim 10^{11}\,\Msun$).
\item Earth orbit altitude: $r_{\mathrm{cross}} \sim 13{,}000$--$22{,}000$~km
  (W3i3).
\item Cosmological void scale: $R_{\mathrm{void}} \sim 100$~Mpc.
\end{itemize}

MAGIS-100 operates at the Earth's surface where $a_N/\azero \sim 10^{11}$,
deep in the Newtonian regime. The $\mu$-function is:
\begin{equation}
\mu(10^{11}) = \frac{10^{11}}{1 + 10^{11}} = 1 - 10^{-11},
\end{equation}
i.e., deviations from Newton are at the $10^{-11}$ level.

\textbf{Direct probe of the crossover}: Not possible from the ground.
The crossover altitude (13,000--22,000~km) requires a satellite platform.
The proposed AEDGE (Atomic Experiment for Dark Matter and Gravity
Exploration in Space) would operate in medium Earth orbit and could
probe the transition region.

\textbf{Indirect constraint on $\mu(x)$ shape}: MAGIS-100 measures
the Newtonian-side behaviour of $\mu(x)$ at $x \sim 10^{11}$. High-precision
measurement of $\delta a/a_N = 1/x = \azero/a_N$ tests the \emph{asymptotic}
form of $\mu(x)$ and constrains the approach to unity. If $\mu(x)$ deviates
from $x/(1+x)$ at large $x$ (e.g., if the true $\mu$ approaches 1 more
slowly), this would be detectable as a larger-than-predicted gravity
gradient anomaly.

\textbf{Connection to T4}: The ISW integral depends on $\mu(x)$ at
$x \sim 0.01$--$1$ (deep-MOND to crossover). MAGIS/AION constrain
$\mu(x)$ only at $x \gg 1$. The connection is through the theoretical
model: if the simple-$\mu$ function $x/(1+x)$ is confirmed at high $x$
by MAGIS/AION, this increases confidence in extrapolating to low $x$,
but does not \emph{prove} the low-$x$ behaviour.

\verdict{The MAGIS/AION connection to T4 is model-validation, not
direct measurement. Resolving T4 requires either (a)~a satellite atom
interferometer in the crossover region, (b)~detailed void-by-void ISW
integration (theoretical), or (c)~improved wide-binary and rotation-curve
constraints on $\mu(x)$ at $x \sim 1$.}
\end{finding}

\subsection{The ISW Integral and $\mu$-Function Sensitivity}

For completeness, the DFD ISW temperature decrement for a photon
traversing a void of radius $R_v$ and central density contrast $\delta_v$:
\begin{equation}
\frac{\Delta T}{T}\bigg|_{\mathrm{ISW}}^{\mathrm{DFD}}
= -\frac{2}{c^3} \int_0^{R_v} \dot{\Phi}_{\mathrm{DFD}}(r, t)\,dr,
\end{equation}
where $\dot{\Phi}_{\mathrm{DFD}}$ includes the $\mu$-function:
\begin{equation}
\dot{\Phi}_{\mathrm{DFD}} = -\frac{c^2}{2}\dot{\psi}
= -\frac{c^2}{2}\frac{\partial}{\partial t}\left[
\mu^{-1}\left(\frac{|\nabla\psi|}{\astar}\right)
\cdot \frac{|\nabla\psi|}{\astar} \cdot \astar
\right].
\end{equation}

The sensitivity of the ISW integral to the $\mu$-function shape at
the crossover:
\begin{equation}
\frac{\partial(\Delta T/T)}{\partial \mu(x=1)}
\propto \frac{\delta_v R_v}{c} \cdot \frac{1}{\mu(1)^2} \cdot H_0,
\end{equation}
where $\mu(1) = 1/2$ for simple-$\mu$. A 10\% change in $\mu(1)$
produces a $\sim$40\% change in the ISW signal. This extreme sensitivity
is why the T4 tension is so difficult to resolve without precise knowledge
of $\mu(x)$ in the crossover regime.

%=============================================================================
\section{P2 SQUID as GGN Pre-Filter: Status Update}
\label{sec:p2-ggn}
%=============================================================================

The P2 SQUID pre-filter synergy (inherited) uses a P2 resonant detector
to monitor gravity gradient noise (GGN) and subtract it from the P6
atom interferometer signal in real time. The GGN noise floor:
\begin{equation}
S_h^{(\mathrm{GGN}),1/2}(f) = \frac{2.87\ee{-21}}{f^{4.3}}\sqrtHz
\end{equation}
dominates the P6 noise budget below $\sim 0.01$~Hz.

\confirmed{The P2 SQUID GGN pre-filter concept remains valid. No new
constraints from W4 iteration. The key requirement is a P2 passive
detector with sensitivity $|\delta\psi| < 10^{-16}$ at $\sim 0.01$~Hz,
which is within the projected capability of a cryogenic toroidal
re-entrant cavity (Nb$_3$Sn, $Q > 10^{11}$).}

%=============================================================================
\section{Derivation Chains and Cross-References}
\label{sec:chains}
%=============================================================================

\subsection{Finding 1 Derivation Chain}

\begin{enumerate}[nosep]
\item DFD field equation: $\nabla\cdot[\mu(|\nabla\psi|/\astar)\nabla\psi]
  = -(8\pi G/c^2)\rho$ (Section~2, formalism).
\item Simple-$\mu$: $\mu(x) = x/(1+x)$ (Appendix A, $S^3$ microsector).
\item Correction at surface: $\delta a/g = \azero/g = 1.14\ee{-11}$
  (Eq.~\ref{eq:delta-a-surface}).
\item Gradient over 100~m: $\delta(\delta a) = 3.6\ee{-16}$~m/s$^2$
  (Eq.~\ref{eq:gradient-numerical}).
\item MAGIS-100 sensitivity: $6\ee{-17}\,g/\sqrt{\mathrm{Hz}}$
  (MAGIS-100 proposal, 1000$\hbar k$ LMT).
\item SNR (1 year): $\sim 540$ (corrected from $\sim 4200$).
\item Solar modulation: below sensitivity by $\sim 10^4$.
\end{enumerate}

\subsection{Finding 2 Derivation Chain}

\begin{enumerate}[nosep]
\item W3i4 noise budget assumes SQL shot-noise-limited atom interferometry.
\item AION prototype (arXiv:2504.09158) demonstrates SQL with ${}^{87}$Sr
  clock transition.
\item Laser phase noise immunity confirmed experimentally.
\item $\Rightarrow$ W3i4 $h_1 = 2.94\ee{-15}\sqrtHz$ validated.
\item AION-10 TDR (arXiv:2508.03491) provides engineering validation
  of 10~m baseline system.
\end{enumerate}

\subsection{Finding 3 Derivation Chain}

\begin{enumerate}[nosep]
\item W3i7 comparison table: DFD P6 loses $10^6$--$10^{10}\times$ on
  raw strain vs all dedicated detectors.
\item January 2026: MAGIS-100 laser lab completed (Fermilab news).
\item August 2025: AION-10 TDR published.
\item ZAIGA Phase 1: 300~m tunnel under construction.
\item No funded decihertz detector achieves $h < 10^{-20}$ before $\sim$2040.
\item DFD P6 niche: scalar signals on existing atom interferometer
  infrastructure.
\end{enumerate}

\subsection{Finding 4 Derivation Chain}

\begin{enumerate}[nosep]
\item LPF noise floor: $5.2$~fm/s$^2/\sqrt{\mathrm{Hz}}$ (PRL 120, 061101).
\item DFD $\dpsiem$ gradient: $\sim 10^{-43}$~m$^{-2}$ (Eq.~\ref{eq:lpf-dpsi} ff.).
\item Resulting acceleration: $4.6\ee{-28}$~m/s$^2$.
\item Ratio to LPF floor: $\sim 10^{-13}$ ($\equiv$ 22 orders below
  strain-equivalent floor, consistent with W3i6).
\item No new constraint beyond SQMS $\lambdaone < 1.56\ee{-9}$.
\end{enumerate}

\subsection{Finding 5 Derivation Chain}

\begin{enumerate}[nosep]
\item T4 tension: 5--8$\times$ ISW overprediction (W3i3--W3i7).
\item MOND crossover altitude: 13,000--22,000~km (W3i3).
\item MAGIS-100 operates at $a_N/\azero \sim 10^{11}$ ($x \gg 1$).
\item ISW integral sensitive to $\mu(x)$ at $x \sim 0.01$--1.
\item MAGIS/AION constrain $\mu$ at $x \gg 1$ only.
\item Connection is model-validation, not direct measurement.
\item Resolution requires satellite (AEDGE) or void-by-void integration.
\end{enumerate}

%=============================================================================
\section{Summary of Corrections and New Results}
\label{sec:summary}
%=============================================================================

\begin{table}[H]
\centering
\caption{W4i10 P6 corrections and new results.}
\label{tab:corrections}
\begin{tabular}{@{}p{4cm}p{3cm}p{5cm}@{}}
\toprule
Item & Previous Value & W4i10 Result \\
\midrule
MAGIS-100 SNR (1 yr) & $\sim 4200$ (W3i4) &
  $\sim 540$ (recalculated; systematic gravity survey errors likely
  dominate) \\[4pt]
AION SQL validation & Assumed (W3i4) &
  Experimentally confirmed (arXiv:2504.09158) \\[4pt]
MAGIS-100 status & Under construction & Laser lab completed Jan 2026;
  commissioning $\sim$2028 \\[4pt]
AION-10 status & Proposed & TDR published Aug 2025; construction at
  Oxford \\[4pt]
LPF constraint on DFD & Not evaluated & $\dpsiem$ signal 22 orders
  below LPF floor; no new constraint \\[4pt]
T4/MOND/MAGIS connection & Not evaluated & Indirect: model-validation
  only; crossover requires satellite \\[4pt]
P6 advantage band & 0.04--0.11~Hz (W3i7) & Confirmed, but for scalar
  signals only \\
\bottomrule
\end{tabular}
\end{table}

\begin{tcolorbox}[colback=green!5!white, colframe=green!60!black,
  title=\textbf{NET ASSESSMENT}]
DFD P6 remains classified as \textbf{NICHE} (unchanged from corpus).
The experimental landscape is advancing rapidly: AION has demonstrated
SQL performance, MAGIS-100 laser lab is complete, and AION-10 is under
construction. The decihertz gap will persist through $\sim$2040. DFD P6's
unique value is \emph{not} raw GW sensitivity (where it loses by $>10^6\times$)
but access to scalar-mode signals on existing atom interferometer infrastructure.
The MAGIS-100 SNR for the DFD $\mu$-function gradient signal is corrected
downward to $\sim$540 (from $\sim$4200), and the solar modulation signature
is below sensitivity. The most productive near-term strategy is to analyse
MAGIS-100 commissioning data (expected $\sim$2028) for the DC gravity gradient
anomaly predicted by DFD.
\end{tcolorbox}

%=============================================================================
% REFERENCES (informal --- full bibliography maintained in master corpus)
%=============================================================================

\section*{References and Sources}

\begin{enumerate}[nosep]
\item Fermilab news, ``Fermilab completes laser lab construction for world's
  largest vertical atom interferometer,'' January 2026.\\
  \url{https://news.fnal.gov/2026/01/fermilab-completes-laser-lab-construction-for-worlds-largest-vertical-atom-interferometer/}

\item AION Collaboration, ``A Prototype Atom Interferometer to Detect Dark
  Matter and Gravitational Waves,'' arXiv:2504.09158 (April 2025).\\
  \url{https://arxiv.org/abs/2504.09158}

\item AION Collaboration, ``AION-10: Technical Design Report for a 10m Atom
  Interferometer in Oxford,'' arXiv:2508.03491 (August 2025).\\
  \url{https://arxiv.org/abs/2508.03491}

\item Armano et al., ``Beyond the Required LISA Free-Fall Performance: New
  LISA Pathfinder Results down to 20~$\mu$Hz,'' Phys.\ Rev.\ Lett.\ 120,
  061101 (2018).\\
  \url{https://journals.aps.org/prl/abstract/10.1103/PhysRevLett.120.061101}

\item Armano et al., ``Sub-Femto-$g$ Free Fall for Space-Based Gravitational
  Wave Observatories: LISA Pathfinder Results,'' Phys.\ Rev.\ Lett.\ 116,
  231101 (2016).\\
  \url{https://link.aps.org/doi/10.1103/PhysRevLett.116.231101}

\item MAGIS-100 Collaboration, ``Matter-wave Atomic Gradiometer Interferometric
  Sensor (MAGIS-100),'' arXiv:2104.02835.\\
  \url{https://arxiv.org/abs/2104.02835}

\item Zhan et al., ``ZAIGA: Zhaoshan Long-baseline Atom Interferometer
  Gravitation Antenna,'' Int.\ J.\ Mod.\ Phys.\ D 29, 1940005 (2020).\\
  \url{https://arxiv.org/abs/1903.09288}

\item DFD Master Findings Corpus, March 2026 (internal).

\item W3i4\_P6\_update.tex, W3i5\_P6\_P7\_update.tex,
  W3i6\_P6\_P7\_P9\_update.tex, W3i7\_P6\_P8\_P9\_update.tex (internal).
\end{enumerate}

\end{document}
